\subsection{Адресация: где какой светодиод?}

Мы выяснили, что в буфере 29 ячеек (\texttt{0..28}). На плате квадрокоптера вы видите светодиодную линейку из 29 элементов.

Как понять, какой индекс в коде соответствует какой лампочке в реальности?

\begin{taskbox}[Исследовательская задача]
Напишите программу-исследователь, которая поможет составить «карту» светодиодов.
\begin{enumerate}
    \item Включите светодиод с индексом \texttt{0} красным цветом.
    \item Посмотрите на дрон. Какой светодиод загорелся? (Например, «Передний левый»).
    \item Повторите для индексов \texttt{1}, \texttt{2}, \texttt{3}, \dots, \texttt{28}.
    \item Запишите результаты в тетрадь.
\end{enumerate}
\end{taskbox}

\begin{center}
\begin{tikzpicture}[scale=1.2]
    % Drone Body (Schematic)
    \draw[fill=gray!10, thick, rounded corners] (-1.5,-1.5) rectangle (1.5,1.5);
    \node at (0,0.5) {\textbf{Pioneer}};
    \node[font=\footnotesize] at (0,0) {(Вид сверху)};

    % Direction Arrow
    \draw[->, ultra thick, GeoskanBlue] (0, 1.6) -- (0, 2.8) node[above, font=\footnotesize\bfseries] {Вперёд};
    
    % Arms
    \draw[thick, double] (1.5,1.5) -- (2.5,2.5);
    \draw[thick, double] (1.5,-1.5) -- (2.5,-2.5);
    \draw[thick, double] (-1.5,1.5) -- (-2.5,2.5);
    \draw[thick, double] (-1.5,-1.5) -- (-2.5,-2.5);
    
    % Motors
    \draw[fill=black!20] (2.5,2.5) circle (0.6) node {\footnotesize M1};
    \draw[fill=black!20] (2.5,-2.5) circle (0.6) node {\footnotesize M2};
    \draw[fill=black!20] (-2.5,2.5) circle (0.6) node {\footnotesize M4};
    \draw[fill=black!20] (-2.5,-2.5) circle (0.6) node {\footnotesize M3};
    
    % LED placeholders (Base)
    \node[circle, draw=red, fill=white, inner sep=3pt, label=left:\textbf{?}] (L1) at (-1.2, 1.2) {};
    \node[circle, draw=red, fill=white, inner sep=3pt, label=right:\textbf{?}] (L2) at (1.2, 1.2) {};
    \node[circle, draw=red, fill=white, inner sep=3pt, label=right:\textbf{?}] (L3) at (1.2, -1.2) {};
    \node[circle, draw=red, fill=white, inner sep=3pt, label=left:\textbf{?}] (L4) at (-1.2, -1.2) {};
    
    \node[red, align=center, below=0.5cm] at (0,-1.5) {Светодиодная линейка\\(индексы 0-28)};

\end{tikzpicture}
\end{center}

\begin{center}
\renewcommand{\arraystretch}{1.5}
\begin{tabular}{|c|p{8cm}|}
\hline
\rowcolor{GeoskanBlue!20} \textbf{Индекс} & \textbf{Расположение (Например: Передний левый, центр, хвост)} \\
\hline
0 & \dotfill \\
\hline
1 & \dotfill \\
\hline
2 & \dotfill \\
\hline
3 & \dotfill \\
\hline
4 & \dotfill \\
\hline
5 & \dotfill \\
\hline
\dots & \dotfill \\
\hline
28 & \dotfill \\
\hline
\end{tabular}
\end{center}

\begin{codebox}[Скрипт-исследователь]
\begin{codefile}{lua/examples/05_leds/leds_addressing_research.lua}
\end{codefile}
\end{codebox}

Это упражнение учит главному принципу инженера: \textbf{если документации недостаточно — проведи эксперимент!} В листинге функция \verb|clear()| по индексу \verb|0..28| выключает все светодиоды, а затем по команде \verb|leds:set(index, ...)| зажигается выбранный номер. Основной принцип — сопоставить индекс из кода с реальным светодиодом на линейке.
